\section{Geometric thresholds and curvature data}
\label{sec:geometric-results}
The location of a shortest-flight threshold is metric information. Its
probability also records the transverse instability of the reflected
segment. These data coincide in a one-parameter circular family, since
the gap already determines the radius. They separate under deformations
of the scatterer. We prove a threshold law on an open set of noncircular
geometries, including the competition between distinct shortest pairs,
and show that its amplitudes determine curvature information which all
threshold locations leave undetermined.

Fix the triangular lattice $\Lambda$ and a reference radius
$R_*\in(0,1/2)$. A scatterer $D_\xi$ is described by its support function
$h_\xi(\theta)$, where $\xi$ ranges over a compact subset $K_0$ of a
finite-dimensional parameter domain. The function $(\xi,\theta)\mapsto
h_\xi(\theta)$ is $C^\infty$ on a neighborhood of $K_0\times S^1$.
There is a $C^4$ neighborhood of the constant $R_*$ in which the
following conclusions hold. No reflection or rotational symmetry is
assumed. Uniform constants for a prescribed derivative order may depend
on the corresponding higher norms of this fixed smooth family; they are
not uniform over unbounded higher derivatives in a $C^4$ ball.

Write $\mathcal E=\{e\in\Lambda:|e|=1\}$, so $|\mathcal E|=6$.
The unique closest segment from $D_\xi$ to $D_\xi+e$ has length
$g_e(\xi)$ and endpoint curvatures $\kappa_{e,0}(\xi)$ and
$\kappa_{e,1}(\xi)$. All these functions are smooth. Put
\begin{equation}\label{eq:g-general-data}
 \begin{split}
 A(\xi)&=\sqrt3/2-|D_\xi|,\qquad g_*(\xi)=\min_{e\in\mathcal E}g_e(\xi),\\
 c_{e,a}&=1+g_e\kappa_{e,a}\quad(a=0,1),\qquad
 c_e=\sqrt{c_{e,0}c_{e,1}},\qquad \gamma_e=\arcosh c_e.
 \end{split}
\end{equation}
Reversing $e$ exchanges the curvatures and leaves $g_e,c_e,\gamma_e$
unchanged. The full preparation is normalized unit-speed phase volume
on $(\R^2/\Lambda)\setminus D_\xi$. We write $N_t$ for the actual
collision count under that preparation.

\begin{theorem}[Stability and competition of collision thresholds]
\label{thm:g-stability}
For all such compact smooth families there is a collar
$\varepsilon_0>0$, independent of every integer $j\ge1$, such that,
if $0\le t-jg_*<\varepsilon_0$, then
\begin{equation}\label{eq:g-competition}
 \begin{split}
 \Pr_\xi(N_t\ge j+1)&=\Pr_\xi(N_t=j+1)\\
 &=\sum_{e\in\mathcal E}
   \frac{d_{j,e,+}^{\,2}}{2A(\xi)\sinh(j\gamma_e(\xi))}
   \bigl[1+d_{j,e,+}H_{j,e}(\xi,d_{j,e,+})\bigr],\\
 d_{j,e}&=t-jg_e(\xi).
 \end{split}
\end{equation}
Here $d_+=\max(d,0)$, and every fixed mixed derivative of $H_{j,e}$
in $\xi$ and its nonnegative offset is bounded independently of $j$.
The functions have smooth right extensions at offset zero. For
$t\le jg_*$ the probability is zero. The collar is smaller than
$\min_{K_0}g_*/4$.

More precisely, the $e$ summand is the probability of the oriented
alternating channel. At its own window $t=jg_e+d$, $0<d<\varepsilon_0$,
that channel has the same formula, whether or not $e$ minimizes $g_e$.
In this last statement the channel specifies the facing contact arcs
of all $j+1$ impacts, not the entire count event at that window.
\end{theorem}

The formula uses the six smooth lengths separately, not derivatives of
the possibly nonsmooth minimum $g_*$. It therefore retains changes of
the minimizing pair. It also keeps the exponentially small coefficient
relative to its own error. For a circle, $c_e=(1-R)/R$ and
$1/(2A)=\lambda_R/(4R)$; summing its six equal terms gives
Theorem~\ref{thm:unweighted} without a change of preparation or horizon.
The complete circular proof is given below, followed by the geometric
proof in Sections~\ref{sec:g-localization}--\ref{sec:g-integration}.

Let $x_{e,0}(\xi),x_{e,1}(\xi)$ be the normal outgoing states, and
let $f_\xi$ be a smooth additive collision mark in their neighborhoods.
Set $a_{e,b}=f_\xi(x_{e,b})$, $\bar f_e=(a_{e,0}+a_{e,1})/2$, and
$w_{j,e}=\sum_{i=0}^j a_{e,i\bmod2}$.

\begin{theorem}[Geometric marked response and its exact source domain]
\label{thm:g-marked}
At the channel's own threshold the marked probability is
\begin{equation}\label{eq:g-marked}
 Z_{j,e}(\xi,q,d)=
 \frac{d^2e^{q w_{j,e}(\xi)}}{2A(\xi)\sinh(j\gamma_e(\xi))}
       [1+dH_{j,e}^{f}(\xi,q,d)].
\end{equation}
Every fixed mixed derivative of the normalized remainder is bounded
uniformly in $j$ on compact parameter and source sets. The complete
marked maximal-count event in the collar of Theorem~\ref{thm:g-stability}
is the sum of these expressions at the positive offsets $d_{j,e}$.
For fixed $\xi$ and real $q$, and for any sufficiently small fixed
$d\ge0$, the series of normalized channel probabilities
$\sum_{j,e}Z_{j,e}(\xi,q,d)/d^2$ converges if and only if
\begin{equation}\label{eq:g-source-domain}
               \max_{e\in\mathcal E}\{q\bar f_e(\xi)-\gamma_e(\xi)\}<0.
\end{equation}
At $d=0$ use right limits. On compact subsets of this open source
domain every fixed mixed derivative can be summed term by term.
The same assertion holds for finitely many real sources, replacing
$q\bar f_e$ by their scalar product.
\end{theorem}

The series concerns channel thresholds in distinct windows. It is not
an unrestricted long-time pressure. Its exact source boundary is useful
precisely because the contributing trajectories have been identified.
The conservative circular source condition of Theorem~\ref{thm:marked}
remains valid and is a sufficient subset of \eqref{eq:g-source-domain}.

Our inverse consequence is not a new interpretation of the radius in a
circle. In Section~\ref{sec:g-inverse} we construct a three-parameter
family of noncircular scatterers with the same shortest gaps at all six
orientations. Their unlabelled count amplitudes determine the unordered
triple of contact curvatures, with the free area also recovered rather
than assumed. In the symmetric one-parameter subfamily, just two
amplitudes suffice. Section~\ref{sec:g-record-response} specifies a
supremum metric for complete records and proves differentiated estimates
for smooth tests and a transverse moving cut. Those are distinct from
convergence for bounded discontinuous decoders.

\paragraph{Organization and notation.}
Sections~\ref{sec:results}--\ref{sec:consequences} treat the circular
specialization, with radius $R$ and the notation introduced there.
Sections~\ref{sec:g-localization}--\ref{sec:g-record-response} prove
the geometric statements for $D_\xi$. The circular convention
$g_R=1-2R$ is not imposed on the latter family. The final section
compares the two threshold regimes with other dynamical limit and
inverse problems.
